% Generated draft: requires mathematical and source review.
\documentclass[11pt]{article}
\usepackage[margin=25mm]{geometry}
\usepackage{fontspec}
\setmainfont{Latin Modern Roman}
\usepackage{amsmath,amssymb,mathtools,unicode-math}
\setmathfont{Latin Modern Math}
\usepackage{enumitem,needspace,xurl,hyperref}
\setlength{\parindent}{0pt}
\setlength{\parskip}{5pt}
\setlength{\emergencystretch}{3em}
\newcommand{\N}{\mathbb{N}}
\newcommand{\Z}{\mathbb{Z}}
\newcommand{\Q}{\mathbb{Q}}
\newcommand{\R}{\mathbb{R}}
\newcommand{\C}{\mathbb{C}}
\newcommand{\F}{\mathbb{F}}
\newcommand{\A}{\mathbb{A}}
\newcommand{\PP}{\mathbb{P}}
\newcommand{\E}{\mathbb{E}}
\newcommand{\eps}{\varepsilon}
\newcommand{\dd}{\,\mathrm{d}}
\begin{document}
\section*{UnsolvedMath ID: 30003059}
\textbf{Source problem code:} catalogue entry
\subsection{Definitions and mathematical statement}
The supplied question asks whether separable equivalence preserves the Castelnuovo--Mumford regularity of Hochschild cohomology for symmetric \(k\)-algebras. It does not specify the hypotheses on \(k\), the finiteness conditions on the algebras or their Hochschild cohomology, or the precise regularity convention. The following gives a precise conditional formulation using finite-dimensional symmetric algebras over a field and Noetherian Hochschild cohomology; these restrictions require confirmation against the source and are not asserted to exhaust its intended generality.

Let \(k\) be a field. A finite-dimensional unital associative \(k\)-algebra \(A\) is symmetric if
\[
A\cong \operatorname{Hom}_k(A,k)
\]
as \(A\)-\(A\)-bimodules, where \((a\cdot f\cdot b)(x)=f(bxa)\). For such an algebra, put \(A^e=A\otimes_k A^{\mathrm{op}}\) and define its Hochschild cohomology ring, with its cohomological grading and cup product, by
\[
\operatorname{HH}^*(A)=\bigoplus_{n\geq 0}\operatorname{Ext}_{A^e}^{n}(A,A).
\]

For this formulation, two finite-dimensional \(k\)-algebras \(A\) and \(B\) are separably equivalent if there exist bimodules \({}_A M_B\) and \({}_B N_A\), finitely generated projective on each side, and bimodule decompositions
\[
M\otimes_B N\cong A\oplus P,
\qquad
N\otimes_A M\cong B\oplus Q,
\]
where \(P\) is an \(A\)-\(A\)-bimodule and \(Q\) is a \(B\)-\(B\)-bimodule. This definition imposes no additional projectivity condition on \(P\) or \(Q\).

For a Noetherian nonnegatively graded, graded-commutative \(k\)-algebra \(R=\bigoplus_{n\geq0}R_n\) with finite-dimensional graded components, let \(R_+=\bigoplus_{n>0}R_n\). Write \(H^i_{R_+}(R)\) for graded local cohomology, obtained by deriving the torsion functor
\[
\Gamma_{R_+}(X)=\{x\in X:(R_+)^m x=0\text{ for some integer }m\geq1\}
\]
in the category of graded \(R\)-modules. Fix the convention
\[
\operatorname{reg}(R)=
\sup\{j+i:i\geq0,\ j\in\mathbb Z,\ H^i_{R_+}(R)_j\neq0\},
\]
with values in \(\mathbb Z\cup\{-\infty,+\infty\}\) and \(\sup\varnothing=-\infty\).

\textbf{Conditional precise question.} For every field \(k\) and every pair of finite-dimensional symmetric \(k\)-algebras \(A,B\) such that \(\operatorname{HH}^*(A)\) and \(\operatorname{HH}^*(B)\) are Noetherian, does separable equivalence imply
\[
\operatorname{reg}\bigl(\operatorname{HH}^*(A)\bigr)
=
\operatorname{reg}\bigl(\operatorname{HH}^*(B)\bigr)?
\]
The intended treatment of algebras outside these finiteness hypotheses, and whether this regularity convention and definition of separable equivalence match the workshop question, remain unresolved from the supplied record. Its reference to preservation ``in general'' must not be read as restricting the question to group algebras or block algebras.
\subsection{Short English statement}
The question asks whether separably equivalent symmetric \(k\)-algebras have Hochschild cohomology rings with equal Castelnuovo--Mumford regularity. A precise formulation is given under explicit field, finite-dimensionality, and Noetherian hypotheses; the supplied record does not establish that these are exactly the intended hypotheses. The supplied dated assessment labels the general question open and distinguishes it from zero-regularity results for group and block algebras.
\subsection{Sources}
\begin{itemize}
\item Local UnsolvedMath registry: \texttt{data/unsolvedmath/problems.json}, permanent ID \(30003059\), title \emph{Hochschild Regularity Under Separable Equivalence}, source problem code \texttt{catalogue entry}, in the Oberwolfach Reports Open Problems collection. The supplied record gives the year as 2016 and names no individual proposer.
\item Supplied source citation: \emph{Hochschild Cohomology in Algebra, Geometry, and Topology} (2016), Oberwolfach Reports, pp.~449--506, DOI \texttt{10.4171/owr/2016/10}. \url{https://doi.org/10.4171/owr/2016/10}. The supplied background describes these reports as records of mathematical workshops and their open-problem discussions.
\item Supplied literature source, labelled ``Tate--Hochschild regularity'': \url{https://openaccess.city.ac.uk/id/eprint/1949/1/1211.5999v1.pdf}.
\item Supplied literature source, labelled ``regularity for block algebras'': \url{https://doi.org/10.1017/CBO9781316227343.020}.
\item Literature assessment dated exactly \texttt{2026-08-21}: the supplied curation labels the general invariance question open and states that known zero-regularity theorems concern group and block cases rather than arbitrary symmetric algebras. This is supplied dated curation, not an independently verified literature assessment. Its background explicitly requests independent verification of time-sensitive status and cited claims.
\end{itemize}
\end{document}
